\section{Policy Performance and Surrogate Fidelity}
\label{sec:experiments_policy_performance_surrogate_fidelity}

This section reports the results of training the initial and surrogate policies for \gls{car}, applied to the six environments.
The first part verifies that the initial policies reach a competent level of play and evaluates how faithfully the surrogate policies reproduce their behavior using the \BehaviorReconstructionName.

\subsection{Initial policy training}
\label{ssec:initial_policy_training}

\begin{figure}[t]
    \centering
    \begin{subfigure}[b]{\textwidth}
        \centering
        \captionsetup{skip=2pt}
        \begin{figure}[t]
    \centering
    \resizebox{\textwidth}{!}{
    \begin{tikzpicture}
    \begin{axis}[
        xlabel={Training iteration},
        ylabel={Mean episode return (\% of best reached)},
        grid=major, ymajorgrids=true,
        xmin=0, ymin=0, ymax=105,
        width=\textwidth, height=0.5\textwidth,
        legend style={at={(0.5,-0.22)}, anchor=north, legend columns=3,
            /tikz/every even column/.append style={column sep=0.4cm}},
    ]
        \addplot[color=breakoutColor, mark=\breakoutMark, mark size=\breakoutMarkSize] coordinates {(1,0.1) (211,19.3) (421,19.3) (631,23.6) (841,30.6) (1051,40.7) (1261,52.1) (1471,23.4) (1681,38.3) (1891,58.2) (2101,45.0) (2311,34.8) (2521,55.6) (2731,54.7) (2941,39.5) (3151,68.5) (3361,60.2) (3571,68.8) (3781,42.2) (3991,52.1) (4201,35.7) (4411,45.7) (4621,45.0) (4831,43.5) (5041,53.8) (5050,74.5)};
        \addlegendentry{Breakout}
        \addplot[color=marioBrosColor, mark=\marioBrosMark, mark size=\marioBrosMarkSize] coordinates {(1,0.3) (188,1.3) (375,2.3) (562,6.3) (749,16.5) (936,42.2) (1123,56.5) (1310,68.2) (1497,69.0) (1684,68.8) (1871,73.5) (2058,68.5) (2245,66.9) (2432,61.5) (2619,72.4) (2806,73.5) (2993,70.6) (3180,71.7) (3367,72.3) (3554,72.5) (3741,75.1) (3928,74.9) (4115,80.4) (4302,85.0) (4489,80.4) (4507,84.9)};
        \addlegendentry{Mario Bros}
        \addplot[color=msPacmanColor, mark=\msPacmanMark, mark size=\msPacmanMarkSize] coordinates {(1,0.8) (268,43.8) (535,69.4) (802,87.5) (1069,93.4) (1336,89.3) (1603,95.1) (1870,89.7) (2137,87.2) (2404,88.6) (2671,95.0) (2938,87.7) (3205,92.7) (3472,75.6) (3739,72.9) (4006,74.7) (4273,79.5) (4540,72.3) (4807,73.8) (5074,81.3) (5341,79.6) (5608,69.1) (5875,71.5) (6142,72.7) (6409,70.5) (6425,59.4)};
        \addlegendentry{Ms.\ Pac-Man}
        \addplot[color=pongColor, mark=\pongMark, mark options={line width=1.5pt, scale=1}, mark size=\pongMarkSize] coordinates {(1,0.0) (22,8.9) (43,20.6) (64,47.1) (85,80.8) (106,84.4) (127,87.8) (148,90.1) (169,91.5) (190,92.1) (211,93.0) (232,93.3) (253,94.3) (274,95.3) (295,95.8) (316,96.9) (337,96.4) (358,97.2) (379,97.5) (400,95.6) (421,97.5) (442,98.6) (463,98.7) (484,98.6) (505,99.5) (507,99.3)};
        \addlegendentry{Pong}
        \addplot[color=spaceInvadersColor, mark=\spaceInvadersMark, mark size=\spaceInvadersMarkSize] coordinates {(1,5.3) (192,29.3) (383,42.8) (574,45.6) (765,58.2) (956,51.4) (1147,41.6) (1338,49.8) (1529,56.5) (1720,53.3) (1911,66.3) (2102,66.9) (2293,74.4) (2484,67.1) (2675,66.3) (2866,58.4) (3057,62.3) (3248,72.8) (3439,79.7) (3630,78.2) (3821,81.9) (4012,81.3) (4203,75.6) (4394,81.6) (4585,75.4) (4604,75.7)};
        \addlegendentry{Space Invaders}
        \addplot[color=surroundColor, mark=\surroundMark, mark size=\surroundMarkSize] coordinates {(1,0.0) (298,17.2) (595,24.1) (892,36.6) (1189,33.3) (1486,44.4) (1783,45.6) (2080,60.1) (2377,73.6) (2674,82.8) (2971,87.2) (3268,88.4) (3565,90.0) (3862,92.8) (4159,92.0) (4456,93.6) (4753,95.3) (5050,94.0) (5347,96.4) (5644,95.2) (5941,97.8) (6238,96.0) (6535,96.3) (6832,98.5) (7128,97.9)};
        \addlegendentry{Surround}
    \end{axis}
    \end{tikzpicture}
    }
    \caption{Training of the initial policies with \gls{ppo}, one curve per environment.}
    \label{fig:initial_policy_training}
\end{figure}

        \caption{Training of the initial policies with \gls{ppo}.}
        \label{fig:initial_policy_training}
    \end{subfigure}

    \vspace{2em}

    \begin{subfigure}[b]{\textwidth}
        \centering
        \captionsetup{skip=2pt}
        \begin{figure}[t]
    \centering
    \resizebox{\textwidth}{!}{
    \begin{tikzpicture}
    \begin{axis}[
        ymode=log,
        xlabel={Epoch}, ylabel={Action loss},
        grid=major, ymajorgrids=true,
        xmin=0,
        width=\textwidth, height=0.5\textwidth,
        legend style={at={(0.5,-0.22)}, anchor=north, legend columns=3,
            /tikz/every even column/.append style={column sep=0.4cm}},
    ]
        \addplot[color=breakoutColor, mark=\breakoutMark, mark size=\breakoutMarkSize] coordinates {(1,3.3005) (34,0.7979) (67,0.4972) (100,0.3760) (133,0.3091) (166,0.2642) (199,0.2330) (232,0.2103) (265,0.1923) (298,0.1786) (331,0.1673) (364,0.1578) (397,0.1491) (430,0.1432) (463,0.1361) (496,0.1310) (529,0.1269) (562,0.1226) (595,0.1192) (600,0.1182)};
        \addlegendentry{Breakout}
        \addplot[color=breakoutColor, forget plot] coordinates {(1,3.3093) (34,0.8592) (67,0.5367) (100,0.4015) (133,0.3262) (166,0.2781) (199,0.2438) (232,0.2191) (265,0.2001) (298,0.1841) (331,0.1720) (364,0.1622) (397,0.1531) (430,0.1459) (463,0.1408) (496,0.1356) (529,0.1292) (562,0.1244) (595,0.1218) (600,0.1202)};
        \addplot[color=breakoutColor, forget plot] coordinates {(1,3.3299) (34,0.9590) (67,0.5697) (100,0.4143) (133,0.3293) (166,0.2777) (199,0.2418) (232,0.2167) (265,0.1969) (298,0.1820) (331,0.1700) (364,0.1597) (397,0.1515) (430,0.1439) (463,0.1382) (496,0.1314) (529,0.1275) (562,0.1228) (595,0.1195) (600,0.1184)};

        \addplot[color=marioBrosColor, mark=\marioBrosMark, mark size=\marioBrosMarkSize] coordinates {(1,14.8064) (6,6.6000) (11,4.6932) (16,3.7979) (21,3.2255) (26,2.7950) (31,2.4568) (36,2.1831) (41,1.9586) (46,1.7765) (51,1.6390) (56,1.5198) (61,1.4210) (66,1.3336) (71,1.2585) (76,1.1925) (81,1.1415) (86,1.0899) (91,1.0473) (96,1.0067) (100,0.9770)};
        \addlegendentry{Mario Bros}
        \addplot[color=marioBrosColor, forget plot] coordinates {(1,13.9914) (6,5.7239) (11,4.2091) (16,3.4875) (21,2.9882) (26,2.5980) (31,2.3031) (36,2.0587) (41,1.8517) (46,1.6883) (51,1.5464) (56,1.4302) (61,1.3393) (66,1.2574) (71,1.1879) (76,1.1276) (81,1.0730) (86,1.0292) (91,0.9897) (96,0.9535) (100,0.9291)};
        \addplot[color=marioBrosColor, forget plot] coordinates {(1,14.4070) (6,5.8056) (11,4.3536) (16,3.5892) (21,3.0552) (26,2.6428) (31,2.3269) (36,2.0684) (41,1.8752) (46,1.6995) (51,1.5680) (56,1.4471) (61,1.3526) (66,1.2709) (71,1.1982) (76,1.1348) (81,1.0816) (86,1.0358) (91,0.9926) (96,0.9550) (100,0.9270)};

        \addplot[color=msPacmanColor, mark=\msPacmanMark, mark size=\msPacmanMarkSize] coordinates {(1,27.1299) (12,4.8831) (23,2.4171) (34,1.4646) (45,1.0649) (56,0.8613) (67,0.7298) (78,0.6299) (89,0.5539) (100,0.4961) (111,0.4505) (122,0.4202) (133,0.3914) (144,0.3688) (155,0.3509) (166,0.3339) (177,0.3196) (188,0.3066) (199,0.2941) (200,0.2929)};
        \addlegendentry{Ms.\ Pac-Man}
        \addplot[color=msPacmanColor, forget plot] coordinates {(1,27.4661) (12,4.8214) (23,2.1927) (34,1.5704) (45,1.7119) (56,0.8459) (67,0.6609) (78,0.5625) (89,0.4881) (100,0.4401) (111,0.6438) (122,0.3744) (133,0.3474) (144,0.3299) (155,0.3120) (166,0.2957) (177,0.2860) (188,0.2749) (199,0.3547) (200,0.2608)};
        \addplot[color=msPacmanColor, forget plot] coordinates {(1,27.8892) (12,4.5449) (23,2.2309) (34,1.4869) (45,1.1135) (56,0.8831) (67,0.7563) (78,0.6963) (89,0.5726) (100,0.5595) (111,0.4733) (122,0.4310) (133,0.4060) (144,0.3792) (155,0.3794) (166,0.3522) (177,0.3243) (188,0.3117) (199,0.2987) (200,0.2966)};

        \addplot[color=pongColor, mark=\pongMark, mark options={line width=1.5pt, scale=1}, mark size=\pongMarkSize] coordinates {(1,2.5732) (6,1.3059) (11,0.9484) (16,0.7769) (21,0.6780) (26,0.6095) (31,0.5613) (36,0.5240) (41,0.4939) (46,0.4681) (51,0.4462) (56,0.4273) (61,0.4109) (66,0.3959) (71,0.3829) (76,0.3711) (81,0.3602) (86,0.3497) (91,0.3405) (96,0.3327) (100,0.3266)};
        \addlegendentry{Pong}
        \addplot[color=pongColor, forget plot] coordinates {(1,2.5733) (6,1.3004) (11,0.9744) (16,0.8094) (21,0.7182) (26,0.6519) (31,0.6026) (36,0.5607) (41,0.5261) (46,0.4978) (51,0.4728) (56,0.4510) (61,0.4325) (66,0.4157) (71,0.4007) (76,0.3878) (81,0.3759) (86,0.3657) (91,0.3560) (96,0.3475) (100,0.3415)};
        \addplot[color=pongColor, forget plot] coordinates {(1,2.5732) (6,1.3249) (11,0.9486) (16,0.7818) (21,0.6848) (26,0.6179) (31,0.5676) (36,0.5267) (41,0.4934) (46,0.4651) (51,0.4414) (56,0.4203) (61,0.4020) (66,0.3867) (71,0.3722) (76,0.3598) (81,0.3487) (86,0.3386) (91,0.3294) (96,0.3211) (100,0.3149)};

        \addplot[color=spaceInvadersColor, mark=\spaceInvadersMark, mark size=\spaceInvadersMarkSize] coordinates {(1,2.1788) (12,1.1213) (23,0.9292) (34,0.8281) (45,0.7546) (56,0.6980) (67,0.6508) (78,0.6127) (89,0.5793) (100,0.5531) (111,0.5309) (122,0.5113) (133,0.4946) (144,0.4799) (155,0.4667) (166,0.4549) (177,0.4442) (188,0.4347) (199,0.4264) (200,0.4260)};
        \addlegendentry{Space Invaders}
        \addplot[color=spaceInvadersColor, forget plot] coordinates {(1,2.1801) (12,1.1053) (23,0.9291) (34,0.8288) (45,0.7589) (56,0.7045) (67,0.6616) (78,0.6256) (89,0.5966) (100,0.5717) (111,0.5507) (122,0.5315) (133,0.5150) (144,0.5005) (155,0.4877) (166,0.4762) (177,0.4654) (188,0.4557) (199,0.4471) (200,0.4466)};
        \addplot[color=spaceInvadersColor, forget plot] coordinates {(1,2.1951) (12,1.1148) (23,0.9097) (34,0.8110) (45,0.7429) (56,0.6917) (67,0.6503) (78,0.6153) (89,0.5868) (100,0.5620) (111,0.5397) (122,0.5190) (133,0.5008) (144,0.4857) (155,0.4726) (166,0.4596) (177,0.4486) (188,0.4391) (199,0.4302) (200,0.4298)};

        \addplot[color=surroundColor, mark=\surroundMark, mark size=\surroundMarkSize] coordinates {(1,24.2212) (34,4.7912) (67,2.0912) (100,1.4291) (133,1.1316) (166,0.9598) (199,0.8497) (232,0.7670) (265,0.7080) (298,0.6612) (331,0.6230) (364,0.5919) (397,0.5663) (430,0.5424) (463,0.5225) (496,0.5059) (529,0.4896) (562,0.4765) (595,0.4648) (600,0.4626)};
        \addlegendentry{Surround}
        \addplot[color=surroundColor, forget plot] coordinates {(1,24.0464) (34,2.8642) (67,1.8330) (100,1.3987) (133,1.1628) (166,1.0158) (199,0.9126) (232,0.8353) (265,0.7736) (298,0.7260) (331,0.6858) (364,0.6513) (397,0.6234) (430,0.5987) (463,0.5766) (496,0.5572) (529,0.5407) (562,0.5243) (595,0.5112) (600,0.5088)};
        \addplot[color=surroundColor, forget plot] coordinates {(1,24.7783) (34,4.1007) (67,1.8595) (100,1.3125) (133,1.0589) (166,0.9075) (199,0.8081) (232,0.7370) (265,0.6835) (298,0.6427) (331,0.6093) (364,0.5809) (397,0.5578) (430,0.5374) (463,0.5206) (496,0.5050) (529,0.4904) (562,0.4782) (595,0.4662) (600,0.4658)};
    \end{axis}
    \end{tikzpicture}
    }
    \caption{Training action loss of the surrogate policies.}
    \label{fig:surrogate_policy_training}
\end{figure}

        \caption{Training action loss of the surrogate policies.}
        \label{fig:surrogate_policy_training}
    \end{subfigure}

    \vspace{1em}
    {\hypersetup{pdfborder={0 0 0}}\ref{policyTrainingLegend}}
    \caption{Training of the initial and surrogate policies.}
    \label{fig:policy_training}
\end{figure}


Figure~\ref{fig:initial_policy_training} shows how the initial policies improve during \gls{ppo} training.
For each environment, the mean episode return is normalized so that $0\%$ and $100\%$ correspond to the lowest and highest mean return observed during that run.
In every environment the return increases and then reaches a plateau, at which point training is stopped.
The number of iterations required to reach this plateau varies from a few hundred, for \PongName, to several thousand, for \SurroundName and \MsPacManName.
For the subsequent analysis, the checkpoint with the best performance is retained rather than the final checkpoint.

To confirm that the trained policies play competently, their scores are compared with the reference scores reported by \cite{machado2017revisitingarcadelearningenvironment}.
Each initial policy reaches a score of the same order as these references, so it constitutes a reasonable target for imitation.

\subsection{Surrogate policy training}
\label{ssec:surrogate_policy_training}

Figure~\ref{fig:surrogate_policy_training} shows, for each environment, the training action loss of the three surrogate policies.
In all cases, the loss decreases smoothly, and the three independently trained policies converge to the same value.

We further evaluate the ability of the surrogate policies to replicate the behavior of the initial policies using the \BehaviorReconstructionName (Table~\ref{tab:behavior_reconstruction}).
The surrogate policies reach a \BehaviorReconstructionName of $0.96$ on average over 100 evaluation episodes, ranging from $0.83$ for \BreakoutName to $1.00$ for \MsPacManName and \SurroundName.
They therefore reproduce the decisions of the initial policies closely enough for their latent representations to be analyzed.

\begin{table}[t]
\centering
\resizebox{0.6\linewidth}{!}{
    \begin{tblr}{
        colspec={|c|c|},
        hlines,
        vlines,
        row{1} = {font=\bfseries, halign=c, valign=m}
    }
        Environment     & \textit{Behavior Reconstruction} \\
        \BreakoutName        & 0.83 $\pm$0.03 \\
        \MarioBrosName      & 0.99 $\pm$0.03 \\
        \MsPacManName    & 1.00 $\pm$0.03 \\
        \PongName            & 0.95 $\pm$0.01 \\
        \SpaceInvadersName  & 0.97 $\pm$0.07 \\
        \SurroundName        & 1.00 $\pm$0.01 \\
    \end{tblr}
}
\caption{Evaluation of \BehaviorReconstructionName across the six Atari environments.}
\label{tab:behavior_reconstruction}
\end{table}

